
\documentclass{article}
\usepackage{colortbl}
\usepackage{makecell}
\usepackage{multirow}
\usepackage{supertabular}

\begin{document}

\newcounter{utterance}

\centering \large Interaction Transcript for game `hot\_air\_balloon', experiment `air\_balloon\_survival\_en\_negotiation\_hard', episode 4 with Qwen3{-}32B{-}t0.0.
\vspace{24pt}

{ \footnotesize  \setcounter{utterance}{1}
\setlength{\tabcolsep}{0pt}
\begin{supertabular}{c@{$\;$}|p{.15\linewidth}@{}p{.15\linewidth}p{.15\linewidth}p{.15\linewidth}p{.15\linewidth}p{.15\linewidth}}
    \# & $\;$A & \multicolumn{4}{c}{Game Master} & $\;\:$B\\
    \hline

    \theutterance \stepcounter{utterance}  
    & & \multicolumn{4}{p{0.6\linewidth}}{
        \cellcolor[rgb]{0.9,0.9,0.9}{
            \makecell[{{p{\linewidth}}}]{
                \texttt{\tiny{[P1$\langle$GM]}}
                \texttt{You are participating in a collaborative negotiation game.} \\
\\ 
\texttt{Together with another participant, you must agree on a single set of items that will be kept. Each of you has your own view of how much each item matters to you (importance). You do not know how the other participant values the items. Additionally, you are given the effort each item demands.} \\
\texttt{You may only agree on a set if the total effort of the selected items does not exceed a shared limit:} \\
\\ 
\texttt{LIMIT = 5138} \\
\\ 
\texttt{Here are the individual item effort values:} \\
\\ 
\texttt{Item effort = \{"C17": 903, "B41": 757, "A71": 850, "C00": 746, "B35": 524, "A77": 204, "A03": 946, "A04": 473, "C54": 616, "C01": 855, "A87": 530, "B32": 419, "B86": 960, "B33": 763, "A10": 730\}} \\
\\ 
\texttt{Here is your personal view on the importance of each item:} \\
\\ 
\texttt{Item importance values = \{"C17": 138, "B41": 583, "A71": 868, "C00": 822, "B35": 783, "A77": 65, "A03": 262, "A04": 121, "C54": 508, "C01": 780, "A87": 461, "B32": 484, "B86": 668, "B33": 389, "A10": 808\}} \\
\\ 
\texttt{Goal:} \\
\\ 
\texttt{Your goal is to negotiate a shared set of items that benefits you as much as possible (i.e., maximizes total importance to YOU), while staying within the LIMIT. You are not required to make a PROPOSAL in every message {-} you can simply negotiate as well. All tactics are allowed!} \\
\\ 
\texttt{Interaction Protocol:} \\
\\ 
\texttt{You may only use the following structured formats in a message:} \\
\\ 
\texttt{PROPOSAL: \{'A', 'B', 'C', …\}} \\
\texttt{Propose keeping exactly those items.} \\
\texttt{REFUSE: \{'A', 'B', 'C', …\}} \\
\texttt{Explicitly reject opponent's proposal.} \\
\texttt{ARGUMENT: \{'...'\}} \\
\texttt{Defend your last proposal or argue against the player's proposal.} \\
\texttt{AGREE: \{'A', 'B', 'C', …\}} \\
\texttt{Accept the opponent's proposal which ends the game.} \\
\texttt{STRATEGIC REASONING: \{'...'\}} \\
\texttt{	Describe your strategic reasoning or anticipation explaining your choice of action. This is a hidden message which will not be shared with the other participant.} \\
\\ 
\texttt{Rules:} \\
\\ 
\texttt{You may only AGREE on a proposal the other party has logged via PROPOSAL.} \\
\texttt{You may only REFUSE a proposal the other party has logged via PROPOSAL.} \\
\texttt{Total effort of any PROPOSAL or AGREE set must be ≤ LIMIT.} \\
\texttt{Do NOT reveal your hidden importance scores.} \\
\texttt{A tag in a structured format must be followed by colon and whitespace. The argument must be a python set containing 0 or more strings.} \\
\texttt{So, it must be of the form TAG: \{...\}} \\
\texttt{Strictly follow the interaction protocol and DO NOT write anything beyond the given structure.} \\
\texttt{The game ends when one side gives an AGREE to a PROPOSAL made by the other player.} \\
\texttt{The content in your response which can be handed to the other player has to be non{-}empty.} \\
\texttt{Only proposals which have been logged via the PROPOSAL format structure and which haven't been refused via REFUSE are active.} \\
\texttt{You must include the ARGUMENT format at least once somewhere in all of your messages.} \\
\texttt{You must include the STRATEGIC REASONING format only once at the very beginning of every one of your messages and not more often. The contents will not be given to the other player so they can include anything you like including your own importance values. Here you should reason step by step to come up with you next move.} \\
            }
        }
    }
    & \\ \\

    \theutterance \stepcounter{utterance}  
    & \multicolumn{4}{p{0.6\linewidth}}{
        \cellcolor[rgb]{0.9,0.9,0.9}{
            \makecell[{{p{\linewidth}}}]{
                \texttt{\tiny{[P1$\rangle$GM]}}
                \texttt{<think>} \\
\texttt{Okay, let's see. I need to figure out the best set of items to propose that maximizes my total importance without exceeding the effort limit of 5138. First, I should calculate the total effort of all items to see if they can all fit. Let me add them up.} \\
\\ 
\texttt{Adding the effort values: C17 (903) + B41 (757) + A71 (850) + C00 (746) + B35 (524) + A77 (204) + A03 (946) + A04 (473) + C54 (616) + C01 (855) + A87 (530) + B32 (419) + B86 (960) + B33 (763) + A10 (730). Let me do this step by step.} \\
\\ 
\texttt{Starting with 903 + 757 = 1660. Then +850 = 2510. +746 = 3256. +524 = 3780. +204 = 3984. +946 = 4930. +473 = 5403. Oh, that's already over 5138. So all items together are way over the limit. So I need to select a subset.} \\
\\ 
\texttt{My goal is to maximize my own importance. Let me list the items with their importance and effort:} \\
\\ 
\texttt{C17: 138 (effort 903)} \\
\texttt{B41: 583 (757)} \\
\texttt{A71: 868 (850)} \\
\texttt{C00: 822 (746)} \\
\texttt{B35: 783 (524)} \\
\texttt{A77: 65 (204)} \\
\texttt{A03: 262 (946)} \\
\texttt{A04: 121 (473)} \\
\texttt{C54: 508 (616)} \\
\texttt{C01: 780 (855)} \\
\texttt{A87: 461 (530)} \\
\texttt{B32: 484 (4} \\
            }
        }
    }
    & & \\ \\

    \theutterance \stepcounter{utterance}  
    & & \multicolumn{4}{p{0.6\linewidth}}{
        \cellcolor[rgb]{0.9,0.9,0.9}{
            \makecell[{{p{\linewidth}}}]{
                \texttt{\tiny{[P1$\langle$GM]}}
                \texttt{Your response did not start with the proper strategic reasoning tag at the very beginning of your response. The very first structured format must be of the form STRATEGIC REASONING: \{...\}. Try again.} \\
            }
        }
    }
    & \\ \\

    \theutterance \stepcounter{utterance}  
    & \multicolumn{4}{p{0.6\linewidth}}{
        \cellcolor[rgb]{0.9,0.9,0.9}{
            \makecell[{{p{\linewidth}}}]{
                \texttt{\tiny{[P1$\rangle$GM]}}
                \texttt{<think>} \\
\texttt{Okay, let me start over. The user wants me to make sure that my response begins with the STRATEGIC REASONING tag. In my previous attempt, I forgot to include that. So first, I need to outline my strategic reasoning step by step.} \\
\\ 
\texttt{I need to maximize my total importance while staying under the effort limit of 5138. Let me list all items with their importance and effort. Then, I'll calculate the total effort of all items to see if they can all fit. If not, I need to find a subset that fits and gives the highest importance.} \\
\\ 
\texttt{First, sum all the efforts. Let's see: 903 (C17) + 757 (B41) + 850 (A71) + 746 (C00) + 524 (B35) + 204 (A77) + 946 (A03) + 473 (A04) + 616 (C54) + 855 (C01) + 530 (A87) + 419 (B32) + 960 (B86) + 763 (B33) + 730 (A10). Adding these up step by step:} \\
\\ 
\texttt{Start with 903 + 757 = 1660. Then +850 = 2510. +746 = 3256. +524 = 3780. +204 = 3984. +946 = 4930. +473 = 5403. That's already over 5138. So all items can't be included. Need to remove some.} \\
\\ 
\texttt{Now, I need to select items that give the highest importance per unit of effort. Let's calculate the importance per effort ratio for each item. For example, A71 has 868 importance and 850 effort, so 868/850 ≈ 1.021. B35 has 783/524 ≈ 1.494. B41 is 583/757 ≈ 0} \\
            }
        }
    }
    & & \\ \\

    \theutterance \stepcounter{utterance}  
    & & \multicolumn{4}{p{0.6\linewidth}}{
        \cellcolor[rgb]{0.9,0.9,0.9}{
            \makecell[{{p{\linewidth}}}]{
                \texttt{\tiny{[P1$\langle$GM]}}
                \texttt{Your response did not start with the proper strategic reasoning tag at the very beginning of your response. The very first structured format must be of the form STRATEGIC REASONING: \{...\}. Try again.} \\
            }
        }
    }
    & \\ \\

    \theutterance \stepcounter{utterance}  
    & \multicolumn{4}{p{0.6\linewidth}}{
        \cellcolor[rgb]{0.9,0.9,0.9}{
            \makecell[{{p{\linewidth}}}]{
                \texttt{\tiny{[P1$\rangle$GM]}}
                \texttt{<think>} \\
\texttt{Okay, I need to make sure I start with the STRATEGIC REASONING tag. Let me outline my plan step by step.} \\
\\ 
\texttt{First, I'll calculate the total effort of all items to see if they can all fit within the limit. If not, I need to find a subset that maximizes my importance without exceeding the effort. The total effort is 5403, which is over 5138. So, I need to remove some items.} \\
\\ 
\texttt{Next, I'll determine which items to remove based on their importance and effort. Items with lower importance per effort ratio should be removed first. For example, A77 has low importance (65) and moderate effort (204). Removing it would save 204 effort and lose 65 importance, which is a good trade{-}off.} \\
\\ 
\texttt{Then, I'll check other items. A03 has high effort (946) but low importance (262). Removing it would save 946 effort and lose 262, which is better than A77. Similarly, A04 has 473 effort and 121 importance. Removing A03 and A04 would save 1419 effort and lose 383 importance. That's a significant saving.} \\
\\ 
\texttt{After removing A77, A03, and A04, the total effort would be 5403 {-} 204 {-} 946 {-} 473 = 5403 {-} 1623 = 3780. That's under the limit. Now, I need to check if there are other items to remove to further optimize. Maybe B86 has high effort (960) but moderate importance (668). Removing it would save 960 and lose 668. However, keeping it might be better if the remaining items can fit. Wait, after removing A77, A03, A04, the total is 3780, which is under 5138. So maybe I can keep all others. Let me verify the total importance in this case.} \\
\\ 
\texttt{Calculating the total importance without A77, A03, A04: 138 (C17) + 583 (B41) + 868 (A71)} \\
            }
        }
    }
    & & \\ \\

    \theutterance \stepcounter{utterance}  
    & & & \multicolumn{2}{p{0.3\linewidth}}{
        \cellcolor[rgb]{0.9,0.9,0.9}{
            \makecell[{{p{\linewidth}}}]{
                \texttt{\tiny{[GM$|$GM]}}
                \texttt{abort} \\
            }
        }
    }
    & & \\ \\

    \theutterance \stepcounter{utterance}  
    & & & \multicolumn{2}{p{0.3\linewidth}}{
        \cellcolor[rgb]{0.9,0.9,0.9}{
            \makecell[{{p{\linewidth}}}]{
                \texttt{\tiny{[GM$|$GM]}}
                \texttt{end game} \\
            }
        }
    }
    & & \\ \\

\end{supertabular}
}

\end{document}
